\documentclass[a4paper, 11pt]{article}	

\usepackage[margin=1.75cm]{geometry} 
\usepackage{listings}
\usepackage{fontspec}
\setmainfont{CMU Serif}
\newfontfamily\cyrillicfonttt{CMU Typewriter Text}[Script=Cyrillic]
\usepackage[russian]{babel}
\usepackage{polyglossia}
\setdefaultlanguage{russian}
\setotherlanguage{english}

\usepackage{amsmath,amssymb}
\usepackage{geometry}
\usepackage{graphicx}
\usepackage{enumitem}
\usepackage{hyperref}
\usepackage{xcolor}
\usepackage{listings}
\usepackage{csquotes}

\geometry{margin=25mm}
\hypersetup{
  colorlinks=true,
  linkcolor=blue,
  urlcolor=blue,
  citecolor=blue
}

\lstdefinelanguage{Lisp}{
  morekeywords={defun,lambda,let,let*,cond,if,progn,car,cdr,cons,list,append,
    null,atom,eq,equal,setq,quote,print,princ,terpri,and,or,not},
  sensitive=true,
  morecomment=[l]{;},
  morestring=[b]"
}

\lstset{
  language=Lisp,
  basicstyle=\ttfamily\small,
  breaklines=true,
  columns=fullflexible,
  frame=single,
  rulecolor=\color{black},
  showstringspaces=false,
  keepspaces=true
}

\setlist[itemize]{leftmargin=*, itemsep=0.2em}
\setlist[enumerate]{leftmargin=*, itemsep=0.2em}

\begin{document}

% -------- титульный лист --------
\begin{titlepage}
    \begin{center}
        \vspace*{0cm}

        \LARGE
        \textbf{Московский государственный университет имени М.В. Ломоносова} \\
        \Large
        Факультет вычислительной математики и кибернетики

        \vspace{6cm}

        \Huge
        \textbf{Отчёт о выполненном задании} \\
        \LARGE \textbf{Программирование символьных вычислений}\\
        \Large 
        Вариант №2: «Геометрия» (a, d, e)\\[3.0em]
        
        \vspace{0.5cm}
        
        \begin{flushright}
        Выполнил: \\
        \textit{Огай Владислав Леонидович}\\
        Группа 324, кафедра АЯ ВМК МГУ
        \end{flushright}
        \vfill

        \Large
        Москва, 2025 г. 

    \end{center}
\end{titlepage}
\tableofcontents
\newpage

% -------- цель работы и постановка задачи --------
\section{Цель работы и постановка задачи}

\subsection*{Цель работы}
Изучить базовые функциональные средства языка программирования Лисп,
ориентированного на символьные вычисления и задачи искусственного интеллекта.
Освоить основные приемы обработки сложных структур символьных данных,
разработать и реализовать лисп-программу для определенного варианта
преобразования символьных выражений.

\subsection*{Постановка задачи}
Пусть точки на плоскости заданы с помощью Лисповского списка:
\[
((x_1\ y_1)\ (x_2\ y_2)\dots).
\]

\noindent\textbf{a. Найти все прямые, порождаемые этими точками, предложить удобную для
последующей обработки структуру хранения этих прямых.}
\begin{enumerate}
  \item найти параллельные и перпендикулярные прямые;
  \item найти прямые, содержащие более двух исходных точек;
  \item найти все точки пересечения получившихся прямых.
\end{enumerate}
\noindent\textbf{d. Дан набор точек на плоскости}
\begin{enumerate}
  \item найти две самые удаленные и самые близкие друг к другу точки;
  \item построить окружность таким образом, что все заданные точки попадут внутрь окружности или окажутся на ее границе;
  \item если такое возможно, среди заданных найти точки, линия, проведенная через которые, делит оставшиеся на два равных множества (точки, оказавшиеся по разные стороны от прямой). Если таких наборов несколько – найти все.
\end{enumerate}
\noindent\textbf{e. По двум наборам из трех точек построить окружности}
\begin{enumerate}
  \item определить, пересекаются ли получившиеся окружности, найти эти точки пересечения;
  \item найти координаты центров обеих окружностей, определить, находится ли центр одной из них внутри другой;
  \item если заданные окружности не пересекаются и одна не лежит внутри другой, то построить третью окружность с минимальным радиусом, которая будет касаться двух заданных, и найти эти точки касания.
\end{enumerate}
\newpage

% -------- представление данных --------
\section{Представление исходных данных программы и вид результирующих данных}
Из постановки задачи известно, что входные данные задаются в виде Лисповского списка,
состоящего из подсписков с 2-мя элементами – координатами точек.
Конкретный состав входных данных, как и вид результирующих данных, зависит от решаемой задачи.
В следствие реализации, программа способна работать исключительно с рациональными числами.
\newline

\textit{Общеиспользуемые структуры данных:}
\begin{itemize}
  \item дробь - список из двух чисел: (num den), где num - числитель, den - знаменатель;
  \begin{itemize}
    \item внутри программы каждое число представлено в виде дроби;
    \item для вывода на экран дроби преобразуются в числа с плавающей точкой.
  \end{itemize}
  \item точка - список из двух чисел: $(x\ y)$, где x - x-координата, y - y-координата.
\end{itemize}

\subsection*{Пункт a}
Задается неограниченный по длине список вида:
\[
((x_1\ y_1)\ (x_2\ y_2)\ (x_3\ y_3)\ (x_4\ y_4)\dots).
\]

\textit{Используемые структуры данных:}
\begin{itemize}
  \item прямая - список из коэффициентов канонического уравнения и списка точек: \[(A\ B\ C\ (x_1\ y_1)\ (x_2\ y_2)\ (x_3\ y_3)\dots).\]
\end{itemize}

\textit{Выходные данные зависят от решаемой подзадачи:}
\begin{itemize}
  \item изначально выводятся все исходные точки и найденные прямые вместе с соответствующими им списками точек;
  \item в 1-й подзадаче отдельно выводятся пары параллельных и перпендикулярных прямых;
  \item во 2-й подзадаче выводятся все прямые, содержащие более двух исходных точек, и соответствующие им списки точек;
  \item в 3-й подзадаче выводятся все точки пересечения получившихся прямых;
  \item в случае, когда данные точки не образуют ни одной прямой, выводится сообщение об ошибке, а подзадачи не решаются.
\end{itemize}

\subsection*{Пункт d}
Задается неограниченный по длине список вида:
\[
((x_1\ y_1)\ (x_2\ y_2)\ (x_3\ y_3)\ (x_4\ y_4)\dots).
\]

\textit{Используемые структуры данных:}
\begin{itemize}
  \item прямая - список из коэффициентов канонического уравнения и списка точек: \[(A\ B\ C\ (x_1\ y_1)\ (x_2\ y_2)\ (x_3\ y_3)\dots);\]
  \item окружность - список из списка координат центра и радиуса: $((x_c\ y_c)\ r)$.
\end{itemize}

\textit{Выходные данные зависят от решаемой подзадачи:}
\begin{itemize}
  \item изначально выводятся все исходные точки;
  \item в 1-й подзадаче отдельно выводятся пара ближайших точек и пара самых отдалённых точек в виде списков из координат;
  \item во 2-й подзадаче выводятся центр окружности (точка в виде списка координат) и радиус (число в обычном представлении Лиспа);
  \item в 3-й подзадаче изначально выводились пары точек, образующих прямые (не обязательно разные), делящие остальные точки на два равных множества;
  \begin{itemize}
    \item в связи с изменением постановки задачи, в 3-й подзадаче выводятся полученные прямые и список точек в виде списков из координат, лежащих на соответствующих прямых; 
  \end{itemize}
  \item подзадачи 1 и 3 могут выдать сообщение об ошибке в случае, когда точек недостаточно для решения подзадачи.
\end{itemize}

\subsection*{Пункт e}
Задается список вида $(a\ b)$, где $a$ и $b$ (наборы точек для задания окружностей) имеют вид:
\[
((x_1\ y_1)\ (x_2\ y_2)\ (x_3\ y_3)).
\]

\textit{Используемые структуры данных:}
\begin{itemize}
  \item окружность - список из списка координат центра и радиуса: $((x_c\ y_c)\ r)$.
\end{itemize}

\textit{Выходные данные зависят от решаемой подзадачи:}
\begin{itemize}
  \item изначально выводятся все исходные точки и окружности, построенные по трём точкам из каждого набора входных данных;
  \item в 1-й подзадаче выводятся точки пересечения окружностей либо сообщение о непересечении;
  \item во 2-й подзадаче выводится сообщение о том, лежит ли центр одной окружности внутри другой;
  \item в 3-й подзадаче: если третья окружность существует – выводится окружность и точки касания, иначе – сообщение о невозможности построения;
  \item все подзадачи могут выдать сообщение об ошибке в случае, когда входные данные некорректны (например, точки лежат на одной прямой или есть повторы).
\end{itemize}

\newpage
% -------- ход работы --------
\section{Ход работы}

\subsection*{Пункт a}
\paragraph{Построение всех прямых по набору точек.}
Реализована функция \texttt{get\_lines}, которая перебирает все пары различных точек и строит прямые:
\[
A = y_2-y_1,\qquad B = x_1-x_2,\qquad C = x_2y_1 - x_1y_2.
\]
Чтобы не хранить дубликаты прямых, используется проверка равенства \texttt{eq\_line} через пропорциональность коэффициентов.
Если прямая уже есть в списке, выполняется объединение списков точек (\texttt{merge\_lines}).

\paragraph{Подзадача a1: параллельные и перпендикулярные прямые.}
\begin{itemize}
  \item Параллельность: \texttt{parallel\_line} по условию $A_1B_2 - A_2B_1 = 0$.
  \item Перпендикулярность: \texttt{perpendicular\_line} по условию $A_1A_2 + B_1B_2 = 0$.
\end{itemize}
Функции \texttt{a1\_parallel} и \texttt{a1\_perpendicular} перебирают все пары прямых и возвращают список пар.

\paragraph{Подзадача a2: прямые, содержащие более двух исходных точек.}
Функция \texttt{a2} фильтрует список прямых, оставляя только те, где список \texttt{points} имеет длину $\ge 3$
(проверка \texttt{has\_over\_2\_points}).

\paragraph{Подзадача a3: все точки пересечения прямых.}
Для каждой пары прямых вычисляется определитель
\[
\det = A_1B_2 - A_2B_1.
\]
Если $\det=0$, прямые параллельны или совпадают и единственной точки пересечения нет.
Иначе точка пересечения находится методом Крамера (\texttt{intersect\_lines}):
\[
x = \frac{B_1C_2 - B_2C_1}{\det},\qquad
y = \frac{C_1A_2 - C_2A_1}{\det}.
\]
Результаты уникализируются через \texttt{add\_point}.

\paragraph{Обёртка для пункта a.}
\texttt{geom\_a} нормализует вход (\texttt{normalize\_points}), строит список прямых и печатает:
все прямые, пары параллельных/перпендикулярных, прямые с $>2$ точками и уникальные пересечения.

\subsection*{Пункт d}

\paragraph{Подзадача d1: ближайшая и самая дальняя пары точек.}
Вместо расстояния со sqrt используется \textbf{квадрат расстояния} (чтобы оставаться в рациональной арифметике):
\[
\operatorname{dist}^2(P_1,P_2)=(x_2-x_1)^2+(y_2-y_1)^2.
\]
Функции \texttt{d1\_near} и \texttt{d1\_further} рекурсивно перебирают пары и выбирают пару самых близких и самых дальних точек через
\texttt{min\_dist}/\texttt{max\_dist}.

\paragraph{Подзадача d2: окружность, содержащая все точки.}
Функция \texttt{d2} строит окружность по ограничивающему прямоугольнику:
\[
C=\left(\frac{\min x+\max x}{2},\,\frac{\min y+\max y}{2}\right).
\]
Радиус в реализации берётся \textbf{с запасом}:
\[
r=\max(\max x-\min x,\ \max y-\min y)+1,
\]
что гарантирует покрытие всех точек (минимальность радиуса не требуется).
Используются \texttt{min\_x/max\_x/min\_y/max\_y}.

\paragraph{Подзадача d3: пары точек, чья прямая делит остальные точки поровну.}
Для пары точек $(p_1,p_2)$ строится прямая \texttt{create\_line}. Для каждой оставшейся точки $p$ вычисляется
\[
s = A x + B y + C.
\]
Если $s>0$ - точка по одну сторону, если $s<0$ - по другую; точки с $s=0$ (лежащие на прямой) не учитываются.
Пара $(p_1,p_2)$ подходит, если количество точек по обе стороны одинаково (\texttt{split\_eq}).
Функция \texttt{d3} возвращает список всех таких пар; для печати пары преобразуются в прямые.

\paragraph{Обёртка для пункта d.}
\texttt{geom\_d} нормализует точки и удаляет повторы (\texttt{unique\_points}), затем печатает результаты d1-d3.

\subsection*{Пункт e}

\paragraph{Построение окружности по трём точкам.}
Окружность хранится как \texttt{(($x_c\ y_c$) r)}.
Перед построением выполняются проверки:
\begin{itemize}
  \item \texttt{three\_points\_ok} - три точки заданы и различны;
  \item \texttt{three\_points\_collinear} - точки не лежат на одной прямой.
\end{itemize}
Центр вычисляется формульно (функция \texttt{get\_center\_coordinate}), радиус - как расстояние от центра до одной из точек.
Важное ограничение реализации: квадратный корень берётся функцией \texttt{sqrt\_frac} \textbf{только если он рационален}
(иначе окружность считается некорректной).

\paragraph{Подзадача e1: взаимное расположение и точки пересечения.}
Сначала вычисляется расстояние между центрами $d$ как \texttt{pair\_dist} (с тем же ограничением рациональности корня).
Далее:
\begin{itemize}
  \item при $d=0$ и $r_1=r_2$ - окружности совпадают (строковое сообщение);
  \item при $d<|r_1-r_2|$ - одна внутри другой (сообщение);
  \item при $d>r_1+r_2$ - не пересекаются и не вложены (сообщение);
  \item при $d=r_1+r_2$ или $d=|r_1-r_2|$ - касание, возвращается одна точка;
  \item иначе - две точки пересечения.
\end{itemize}
Точки пересечения вычисляются \texttt{inter\_circles} (классический алгоритм через $(d,a,h)$), при этом $h=\sqrt{r_1^2-a^2}$
также должен быть рациональным, иначе пересечение не возвращается.

\paragraph{Подзадача e2: лежит ли центр одной окружности внутри другой.}
Функция \texttt{e2} сравнивает $d$ с радиусами: центр окружности лежит внутри другой, если $d<r$.

\paragraph{Подзадача e3: третья окружность минимального радиуса, касающаяся двух данных.}
Строится только в случае внешнего раздельного расположения ($d>r_1+r_2$) и рационального $d$:
\[
r_3=\frac{d-r_1-r_2}{2},\qquad
C_3 = C_1 + \frac{r_1+r_3}{d}(C_2-C_1).
\]
Данная окружность лежит на прямой, образованной центрами двух исходных окружностей.
\newline
Реализация: \texttt{third\_circle}. Точки касания получаются как пересечения окружности 3 с окружностями 1 и 2
(\texttt{third\_circle\_inter}); при касании \texttt{inter\_circles} даёт две одинаковые точки, берётся первая.

\paragraph{Обёртка для пункта e.}
\texttt{geom\_e} ожидает 6 точек (две тройки), строит две окружности через \texttt{e0\_try} (с диагностикой ошибок),
после чего запускает и печатает результаты e1, e2, e3.
\newpage

% -------- тесты --------
\section{Тестовые наборы}

\subsection*{Тестовые наборы для пункта a}
\begin{lstlisting}
; 1) Пустой список — недостаточно точек
; ожидается сообщение: ОШИБКА: Недостаточно точек для построения прямой
(geom_a '())

; 2) Две одинаковые точки — прямые построить нельзя (слишком мало различных точек)
; ожидается сообщение: ОШИБКА: Невозможно построить прямые (слишком мало различных точек)
(geom_a '((0 0) (0 0)))

; 3) Две различные точки — одна прямая x - y = 0
; параллельных/ перпендикулярных пар нет, прямых с >2 точками нет, точек пересечения нет
(geom_a '((0 0) (1 1)))

; 4) Три коллинеарные точки — одна прямая x - y = 0, содержащая 3 точки
; подпункт a2 выводит эту прямую, точек пересечения нет
(geom_a '((0 0) (1 1) (2 2)))

; 5) Прямоугольник: (0,0), (2,0), (2,1), (0,1)
; прямые ( в упрощённом виде): y=0, y=1, x=0, x=2, x-2y=0, x+2y-2=0
; параллельные пары: (y=0 || y=1), (x=0 || x=2)
; перпендикулярные пары: y=0 x=0, y=0 x=2, y=1 x=0, y=1 x=2
; точки пересечения:
; (0 1) (1 1/2) (2 1) (2 0) (0 0)
(geom_a '((0 0) (2 0) (2 1) (0 1)))
\end{lstlisting}

\subsection*{Тестовые наборы для пункта d}
\begin{lstlisting}
; 1) Пустой список
; d1: ОШИБКА: недостаточно точек для поиска пар
; d2: ОШИБКА: окружность не построена
; d3: ОШИБКА: недостаточно точек для классификации
(geom_d '())

; 2) Только одна уникальная точка (с учётом дубликатов)
; d1: ОШИБКА: недостаточно точек для поиска пар
; d2: окружность: центр (1 2), радиус 1
; d3: ОШИБКА: недостаточно точек для классификации
(geom_d '((1 2) (1 2)))

; 3) Три точки
; ближайшая пара: (5 8) (-3 3)
; самая дальняя пара: (6 -8) (5 8)
; d2: центр (3/2 0), радиус 17
; d3: разделяющих прямых нет (ничего не выводится после заголовка подпункта d3)
(geom_d '((-3 3) (5 8) (6 -8)))

; 4) Пять точек, классификация возможна (несмотря на нечётное число точек: часть точек лежит на разделяющей прямой)
; ближайшая пара: (8 5) (6 5)
; самая дальняя пара: (8 5) (-2 1)
; d2: центр (3 3), радиус 11
; d3: одна разделяющая прямая (в виде из вывода коэффициентов): -2x + 5y - 9 = 0
; (эквивалентно y = 0.4x + 1.8), на ней в выводе будут точки: (8 5) (3 3) (-2 1)
(geom_d '((-2 1) (3 3) (6 2) (6 5) (8 5)))

; 5) Набор с дубликатом (после unique_points остаётся 4 точки), классификация возможна
; ближайшая пара: (3 -4) (3 -2)
; самая дальняя пара: (3 -4) (2 2)
; d2: центр (4 -1), радиус 7
; d3: три разделяющие прямые (в виде из вывода коэффициентов):
;  4x + 1y - 10 = 0   (y = -4x + 10)   по точкам (3 -2) (2 2)
;  0x + 3y + 6 = 0    (y = -2)         по точкам (6 -2) (3 -2)
;  2x + 0y - 6 = 0    (x = 3)          по точкам (3 -4) (3 -2)
(geom_d '((2 2) (3 -2) (6 -2) (3 -2) (3 -4)))
\end{lstlisting}

\subsection*{Тестовые наборы для пункта e}
\begin{lstlisting}
; 1) Ошибка: 1-я тройка коллинеарна
; ожидается сообщение:
; ОШИБКА: первую окружность построить невозможно: точки лежат на одной прямой!
; ОШИБКА: подзадачи e1, e2 и e3 решить нельзя — есть некорректные окружности!
(geom_e '((0 0) (0 0) (1 0) (5 0) (0 5) (-5 0)))

; 2) Ошибка: 2-я тройка содержит совпадающие точки
; ожидается сообщение:
; ОШИБКА: вторую окружность построить невозможно из-за совпадающих точек!
; ОШИБКА: подзадачи e1, e2 и e3 решить нельзя — есть некорректные окружности!
(geom_e '((5 0) (0 5) (-5 0) (0 0) (0 0) (1 0)))

; 3) Две окружности пересекаются в двух точках
; Окружность 1: центр (1 2), радиус 5
; Окружность 2: центр (7 2), радиус 5
; e1: точки пересечения: (4 -2) и (4 6)
; e2: Ни один из центров не лежит внутри другой окружности
; e3: Окружности не располагаются так, чтобы можно было построить подходящую окружность!
(geom_e '((6 2) (1 7) (-4 2) (12 2) (7 7) (2 2)))

; 4) Внешнее касание (одна точка)
; Окружность 1: центр (0 0), радиус 3
; Окружность 2: центр (5 0), радиус 2
; e1: касание в точке (3 0)
; e2: Ни один из центров не лежит внутри другой окружности
; e3: нельзя построить 3-ю окружность (условие d > r1+r2 не выполняется)
(geom_e '((3 0) (0 3) (-3 0) (7 0) (5 2) (3 0)))

; 5) Окружности совпадают
; обе окружности: центр (0 0), радиус 5
; e1: Окружности совпадают
; e2: Оба центра лежат в других окружностях (так выводит текущая реализация)
; e3: нельзя построить 3-ю окружность
(geom_e '((5 0) (0 5) (-5 0) (0 5) (-5 0) (0 -5)))

; 6) Окружности не пересекаются и не вложены — строится 3-я окружность минимального радиуса
; Окружность 1: центр (0 0), радиус 3
; Окружность 2: центр (10 0), радиус 2
; e1: Окружности не пересекаются и не вложены друг в друга
; e2: Ни один из центров не лежит внутри другой окружности
; e3: третья окружность: центр (11/2 0), радиус 5/2, точки касания: (3 0) и (8 0)
(geom_e '((3 0) (0 3) (-3 0) (12 0) (10 2) (8 0)))

; 7) Невозможно определить взаимное расположение: расстояние между центрами не выражается рационально
; Окружность 1: центр (0 0), радиус 5
; Окружность 2: центр (1 1), радиус 1
; e1: Невозможно определить отношение окружностей: расстояние не выражается рационально
; e2: Невозможно определить: расстояние между центрами не выражается рационально
; e3: Окружности не располагаются так, чтобы можно было построить подходящую окружность!
(geom_e '((5 0) (0 5) (-5 0) (2 1) (1 2) (0 1)))
\end{lstlisting}

\newpage
% -------- дополнения --------
\section{Дополнения к программе}
К описанной выше программе было написано 2 дополнения:
\begin{enumerate}
  \item убрать повторы прямых, удовлетворяющих условию 3-й подзадачи пункта d);
  \item определить по набору точек, каким прямым, образованным данными точками (не считая первой), принадлежит первая точка.
\end{enumerate}

Входными данными для обоих дополнений является список координат точек вида
\[
((x_1\ y_1)\ (x_2\ y_2)\ (x_3\ y_3)\ (x_4\ y_4)\dots).
\]
В программе перед вычислениями точки приводятся к рациональному виду с помощью
\texttt{(normalize\_points points)}, при необходимости дубликаты устраняются функцией
\texttt{(unique\_points points)} (после нормализации).

\subsection*{Дополнение 1: удаление повторов прямых в подпункте d3}
Подпункт d3 формировал список \emph{пар точек}, каждая такая пара задавала разделяющую прямую.
Если на одной и той же разделяющей прямой лежит более двух исходных точек, то эта прямая
может появляться несколько раз (для разных пар коллинеарных точек).

Для устранения повторов при выводе используется преобразование пар в прямые и объединение
совпадающих прямых:
\begin{itemize}
  \item \texttt{print\_classification\_result} вызывает \texttt{get\_class\_res};
  \item \texttt{get\_class\_res} для каждой пары строит прямую через \texttt{(get\_lines pair nil)}
        и добавляет её в список через \texttt{add\_line};
  \item \texttt{add\_line} проверяет равенство прямых (\texttt{eq\_line}) и при совпадении
        объединяет списки точек на прямой (\texttt{merge\_lines}).
\end{itemize}
В результате каждая разделяющая прямая выводится ровно один раз.

\subsection*{Дополнение 2: прямые, содержащие первую точку}
Определена функция \texttt{(on\_the\_line data)}:
\begin{itemize}
  \item первая точка списка \texttt{data} рассматривается как проверяемая точка $p$;
  \item по остальным точкам строятся все прямые: \texttt{(get\_lines (cdr data) nil)};
  \item отбираются только те прямые, для которых выполняется уравнение
        $Ax + By + C = 0$ в точке $p$ (предикат \texttt{check\_line});
  \item результат выводится через \texttt{print\_lines}. Если подходящих прямых нет, функция ничего не выводит.
\end{itemize}

\subsection*{Тестовые наборы для дополнения 1}
Фактически дополнение 1 изменяет только вывод подпункта d3. Соотственно, тесты идентичны оным из пункта d, но с учётом отсутствия повторов прямых.

\subsection*{Тестовые наборы для дополнения 2}
\begin{lstlisting}
; Первая точка (2 2) лежит на двух прямых, построенных по точкам (0 0), (1 1), (2 2), (3 3),
; поэтому должны быть выведена прямые y=x и y=-x+4 ( в формате программы: Ax + By + C = 0)
(on_the_line '((2 2) (0 0) (1 1) (2 2) (3 3) (0 4)))


; Первая точка (0 1) не лежит ни на одной прямой, образованной точками (0 0), (1 0), (1 1),
; вывод отсутствует ( пустой список прямых).
(on_the_line '((0 1) (0 0) (1 0) (1 1)))
\end{lstlisting}

\end{document}
